\chapter{Risks, Requirements \& Constraints}

\section{Constraints \& Requirements}
\begin{itemize}
    \item \textbf{Timeline \& Compute Factor:} The project is strictly constrained by a two-semester timeline. A major requirement is factoring in the substantial time required to train multiple computationally expensive 3DGS models before the change detection evaluation phase can even begin. 
    \item \textbf{Data Availability \& Suitability:} The project relies heavily on acquiring suitable temporal data. A critical constraint is finding data that surveys an identical geographical location with a sufficient time gap to demonstrate structural change.
    \item \textbf{Data Security \& Compliance:} Any data utilized, particularly the AIMS EcoRRAP dataset, must be securely managed and used strictly within the bounds of the established data-sharing or borrowing agreements.
    \item \textbf{Hardware Access:} Training 3DGS models requires continuous access to high-end GPUs. A strict constraint is securing this compute time locally; failure to do so will require pivoting the project budget into cloud computing solutions (e.g., AWS, RunPod) to train the pipeline.
\end{itemize}

\section{Risks and Blockers}
\begin{itemize}
    \item \textbf{Algorithm Failing on Underwater Noise:} If the inherent underwater noise, scattering, and optical attenuation cannot be adequately filtered during the initial reconstruction phase, these errors will propagate directly into the GS-DIFF model, causing the algorithm to falsely identify noise as structural change.
    \item \textbf{Extreme Environmental Change:} Too much change between timestamps could result in completely uncorrelated reconstructions, making temporal alignment and tracking impossible.
    \item \textbf{Moving Features:} Transient features in the dataset (like marine life, moving sand, or shifting shadows) could duplicate on the change map, heavily cluttering the output and masking actual coral growth or decay.
    \item \textbf{Complexity of Coral:} Coral is highly geometrically complex and can be exceedingly difficult to model accurately via Gaussian splatting, especially when measured under noisy underwater conditions.
\end{itemize}
